\chapter{Propensity Score Matching}
\label{cap:psm}
% ── índice ──────────────────────────────────────────────────────
\index{matching por escore de propensão|textbf}
\index{PSM|see{matching por escore de propensão}}
\index{psm@\texttt{psm()}|textbf}
\index{escore de propensão|textbf}
\index{ATT|see{efeito médio do tratamento sobre tratados}}
\index{caliper}
\index{viés de seleção}

% ─────────────────────────────────────────────────────────────────────────────
\section{Seleção em observáveis}

Em muitos estudos observacionais, a atribuição ao tratamento não é aleatória
— indivíduos se selecionam ao tratamento com base em características
observáveis $X$. Se $X$ afeta também o resultado $Y$, a comparação direta
entre tratados e controles produz estimativas viesadas.

O \textbf{Propensity Score Matching} (PSM) explora o teorema de Rosenbaum
e Rubin (1983): se a atribuição ao tratamento é independente dos
resultados potenciais dado $X$ (\emph{ignorabilidade}), então a mesma
independência vale dado o propensity score $p(X) = P(D=1 \mid X)$:

\[
  (Y^0, Y^1) \perp D \mid X
  \quad\Longrightarrow\quad
  (Y^0, Y^1) \perp D \mid p(X)
\]

Isso reduz o problema de dimensão alta (condicionar em $X$) para um
problema unidimensional (condicionar em $\hat{p}$). O PSM:
\begin{enumerate}
  \item Estima $\hat{p}(X_i)$ via logit (score de propensão).
  \item Emparelha cada tratado com o(s) controle(s) mais próximo(s) em $\hat{p}$.
  \item Computa o ATT como diferença de médias no grupo emparelhado.
\end{enumerate}

A condição de suporte comum (\emph{overlap}): $0 < p(X) < 1$ para todo $X$
é essencial — unidades com probabilidade 0 ou 1 de tratamento não têm
contrafactual no outro grupo.

% ─────────────────────────────────────────────────────────────────────────────
\section{Verificação por DGP}

O DGP cria confusão via $x_1$: quem tem $x_1$ alto tem maior probabilidade
de tratamento e também maior resultado em média. O ATT verdadeiro é $\tau = 2$.

\[
  P(D=1 \mid x_1) = 0{,}3 + 0{,}4\,x_1 \in [0{,}30,\; 0{,}70]
  \qquad
  y_i = 1 + 2\,d_i + 3\,x_{1i} + \varepsilon_i
\]

\begin{lstlisting}[caption={DGP com seleção em observáveis}]
seed(42)
input df
id
1
end
for i in 1..=10 { df = append(df, df) }
df |> filter(_n <= 1000)
generate df x1 = runiform()                   // x1 ~ U(0,1)
generate df u  = runiform()                   // ruído de atribuição
generate df d  = (0.3 + 0.4*x1) > u          // P(D=1|x1) ∈ [0.30, 0.70]
generate df e  = rnormal()
generate df y  = 1.0 + 2.0*d + 3.0*x1 + e   // ATT verdadeiro = 2
\end{lstlisting}

\subsection*{OLS ingênuo vs.\ OLS com controles}

\begin{lstlisting}[caption={Comparação OLS sem e com controle de x1}]
let m_ols = ols(y ~ d, df)          // sem controle: viesado
let m_adj = ols(y ~ d + x1, df)    // com controle: não viesado
display m_ols
display m_adj
\end{lstlisting}

\begin{lstlisting}[language={}, basicstyle=\ttfamily\small,
  frame=none, numbers=none, backgroundcolor=\color{codbg},
  xleftmargin=0pt, xrightmargin=0pt, breaklines=false]
── OLS sem controle ──────────────────────────────────
const  2.9135***  (0.0398)    d  2.2215***  (0.0562)
N=1000  R²=0.6105
──────────────────────────────────────────────────────

── OLS com x1 ────────────────────────────────────────
const  1.5289***  (0.0193)    d  1.9762***  (0.0180)
x1     3.0053***  (0.0317)
N=1000  R²=0.9610
──────────────────────────────────────────────────────
\end{lstlisting}

O OLS sem controle superestima $\tau$ em $+0{,}22$ (2.22 vs.\ 2.0) — viés
de variável omitida. Controlando $x_1$ o OLS recupera $\hat\tau = 1{,}98$,
indistinguível de 2.

\subsection*{PSM sem caliper}

\begin{lstlisting}[caption={PSM 1:1 sem caliper}]
let m_psm = psm(y ~ d + x1, df, k=1, boot=200)
display m_psm
\end{lstlisting}

Sem caliper, o matching aceita pares distantes em $\hat{p}$. O balanceamento
pode melhorar pouco e a estimativa pode permanecer viesada se tratados forem
pareados com controles ruins.

\subsection*{PSM com caliper}

O caliper limita a distância máxima admitida no matching. A regra prática
de Cochran e Rubin (1973) é caliper $= 0{,}2 \times \sigma_{\hat{p}}$, o
que corresponde aproximadamente a $0{,}05$ para scores uniformes.

\begin{lstlisting}[caption={PSM com caliper = 0.05}]
let m_cal = psm(y ~ d + x1, df, k=1, caliper=0.05, boot=200)
display m_cal
\end{lstlisting}

Com caliper $= 0{,}05$, a estimativa deve se aproximar do ATT verdadeiro ao
custo de descartar tratados sem par aceitável. O caliper também deve melhorar
o balanceamento ao rejeitar matches distantes no propensity score.\footnote{%
  O aviso \texttt{SMD > 0.10} persiste porque a regra convencional exige
  SMD $\leq 0{,}10$ (Austin 2009). Neste DGP a confusão por $x_1$ é intensa
  ($\beta_{x_1} = 3$) e $N = 1000$; com mais observações o caliper poderia
  ser estreitado para melhorar ainda mais o balanço.}

% ─────────────────────────────────────────────────────────────────────────────
\section{Sintaxe}

\begin{lstlisting}
psm(outcome ~ treatment + cov1 + cov2, df)
psm(outcome ~ treatment + cov1 + cov2, df, k=1, caliper=0.05, replace=false, boot=200)
\end{lstlisting}

A fórmula: \lstinline{outcome ~ treatment + cov1 + cov2 + ...}. O primeiro
regressor após \lstinline{~} é a variável de tratamento; os demais são as
covariáveis usadas no modelo de propensão.

\begin{center}
\begin{tabular}{lll}
\toprule
\textbf{Opção} & \textbf{Padrão} & \textbf{Descrição} \\
\midrule
\texttt{k}        & 1       & número de controles por tratado \\
\texttt{caliper}  & nenhum  & distância máxima no score de propensão \\
\texttt{replace}  & false   & matching com reposição \\
\texttt{boot}     & 200     & réplicas bootstrap para o SE \\
\bottomrule
\end{tabular}
\end{center}

\begin{lstlisting}[caption={Fluxo típico — avaliação de programa}]
load "beneficiarios.csv" as df

// Covariáveis de seleção
let m = psm(renda ~ tratado + idade + educ + regiao, df,
            k=1, caliper=0.05, boot=500)
display m
\end{lstlisting}

\begin{nota}
  O PSM é consistente sob \emph{ignorabilidade} (seleção em observáveis):
  $\{Y^0, Y^1\} \perp D \mid X$. Variáveis omitidas que afetam
  simultaneamente $D$ e $Y$ invalidam a identificação — o mesmo problema
  que torna IV necessário em cap.~\ref{cap:iv}. Ao contrário do DiD
  (cap.~\ref{cap:did}), o PSM não elimina confusores não observados
  constantes no tempo.
\end{nota}

\section{Validação empírica}

O dataset \texttt{jtrain3} do Wooldridge é usado no caso de validação
\texttt{psm_jtrain3} em \texttt{validation/cases/}. O caso estima o ATT de
um programa de treinamento por PSM em \hayashi{} e o compara com uma
implementação de referência em Python. Execute \lstinline{hay validate} para
ver o estado atual.
